% Part: propositional-logic
% Chapter: propositional-logic
% Section: preliminaries

\documentclass[../../../include/open-logic-section]{subfiles}

\begin{document}

\olfileid{pl}{syn}{pre}

\olsection{تمهيدات}

\begin{thm}[\emph{مبدأ الاستقراء على !!{formula}s}]
  \ollabel{thm:induction}
  إذا كانت خاصية ما~$P$ متحققة لجميع ما هو ذري من !!{formula}s، وكانت
  بحيث
  \begin{enumerate}
    \tagitem{prvNot}{تتحقق لـ~$\lnot !A$ كلما تحققت
      لـ~$!A$؛}{}
    \tagitem{prvAnd}{تتحقق لـ~$(!A \land !B)$
      كلما تحققت لـ~$!A$ و~$!B$؛}{}
    \tagitem{prvOr}{تتحقق لـ~$(!A \lor !B)$
      كلما تحققت لـ~$!A$ و~$!B$؛}{}
    \tagitem{prvIf}{تتحقق لـ~$(!A \lif !B)$
      كلما تحققت لـ~$!A$ و~$!B$؛}{}
    \tagitem{prvIff}{تتحقق لـ~$(!A \liff !B)$
      كلما تحققت لـ~$!A$ و~$!B$؛}{}
  \end{enumerate}
  فإن $P$ تتحقق لجميع !!{formula}s.
\end{thm}

\begin{proof}
  لتكن $S$ مجموعة كل ما يتمتع بالخاصية~$P$ من !!{formula}s. من الواضح
  أن $S \subseteq \Frm[L_0]$. تستوفي~$S$ جميع شروط
  \olref[fml]{defn:formulas}: فهي تحتوي على كل ما هو ذري من
  !!{formula}s ومنغلقة تحت !!{operator}s. و~$\Frm[L_0]$ هو
  أصغر صنف من هذا القبيل، ومن ثم $\Frm[L_0] \subseteq S$. وعليه
  $\Frm[L_0] = S$، وكل صيغة تتمتع بالخاصية~$P$.
\end{proof}

\begin{prop}\ollabel{prop:balanced}
   كل !!{formula} في~$\Frm[L_0]$ \emph{متوازنة}، بمعنى أن فيها من
   الأقواس اليسرى بقدر ما فيها من الأقواس اليمنى.
\end{prop}

\begin{prob}
أثبت \olref[pl][syn][pre]{prop:balanced}
\end{prob}

\begin{prop} \ollabel{prop:noinit}
لا يكون أي مقطع ابتدائي حقيقي من !!a{formula} هو نفسه !!a{formula}.
\end{prop}

\begin{prob}
  أثبت \olref[pl][syn][pre]{prop:noinit}
\end{prob}

\begin{prop}[وحدانية القراءة]
لكل !!{formula}~$!A$ في $ \Frm[L_0]$ تحليل واحد بالضبط بوصفها واحدة من
الصور الآتية
\begin{enumerate}
\tagitem{prvFalse}{$\lfalse$.}{}

\tagitem{prvTrue}{$\ltrue$.}{}

\item $\Obj p_n$ لمتغير ما $\Obj p_n \in  \PVar$.

\tagitem{prvNot}{$\lnot !B$ لـ~$!B$ بوصفه !!{formula} ما.}{}

\tagitem{prvAnd}{$(!B \land !C)$ لـ~$!B$ و~$!C$ من !!{formula}s.}{}

\tagitem{prvOr}{$(!B \lor !C)$ لـ~$!B$ و~$!C$ من !!{formula}s.}{}

\tagitem{prvIf}{$(!B \lif !C)$ لـ~$!B$ و~$!C$ من !!{formula}s.}{}

\tagitem{prvIff}{$(!B \liff !C)$ لـ~$!B$ و~$!C$ من !!{formula}s.}{}
\end{enumerate}
وفوق ذلك، هذا التحليل \emph{وحيد}.
\end{prop}

\begin{proof}
بالاستقراء على $!A$. فمثلًا، افترض أن لـ~$!A$ تحليلين متميزين على
أنها $(!B \lif !C)$ و~$(!B' \lif !C')$. عندئذ يجب أن يكون $!B$ و~$!B'$
متطابقين (وإلا لكان أحدهما مقطعًا ابتدائيًا حقيقيًا من الآخر)؛ فإذا كان
تحليلا~$!A$ متميزين فلا بد أن يكون ذلك لأن $!C$ و~$!C'$ تحليلان
متميزان لسلسلة الرموز نفسها، وهذا مستحيل بفرض الاستقراء.
\end{proof}


\begin{defn}[الإحلال المنتظم]
إذا كان $!A$ و~$!B$ من !!{formula}s، وكان $\Obj p_i$ !!{propositional
variable}، فإن $\Subst{!A}{!B}{\Obj p_i}$ يرمز إلى نتيجة
استبدال كل ظهور لـ~$\Obj p_i$ بظهور لـ~$!B$ في~$!A$؛
وبالمثل، يرمز إلى الإحلال المتزامن لـ~$\Obj p_1$، \dots،~$\Obj p_n$ بما
يقابلها من !!{formula}s $!B_1$، \dots،~$!B_n$ بالرمز
$\SSubst{!A}{\subst{!B_1}{\Obj p_1},\dots,\subst{!B_n}{\Obj p_n}}$.
\end{defn}

\begin{prob} لكل واحدة من خمس !!{formula}s أدناه، حدّد ما إذا كان يمكن
 التعبير عن !!{formula} على صورة إحلال \( \Subst{!A}{!B}{\Obj p_i} \)
 حيث تكون \( !A \) هي (i) \( \Obj p_0 \)؛ و(ii) \( ( \lnot \Obj p_0 \land \Obj
 p_1) \)؛ و(iii) \( ( ( \lnot \Obj p_0 \lif \Obj p_1 ) \land \Obj
 p_2 ) \). وفي كل حالة حدّد الإحلال ذا الصلة.
  \begin{enumerate}
    \item \( \Obj p_1 \)
    \item \( ( \lnot \Obj p_0 \land \Obj p_0 ) \)
    \item \( ( ( \Obj p_0 \lor \Obj p_1 ) \land \Obj p_2 )  \)
    \item \( \lnot ( ( \Obj p_0 \lif \Obj p_1 ) \land \Obj p_2 ) \)
    \item \( (( \lnot ( \Obj p_0 \lif \Obj p_1 ) \lif ( \Obj p_0 \lor \Obj p_1 )) \land \lnot ( \Obj p_0 \land \Obj p_1 )) \)
  \end{enumerate}
\end{prob}

\begin{prob}
أعط تعريفًا صارمًا رياضيًا لـ~$\Subst{!A}{!B}{p}$ بطريق
الاستقراء.
\end{prob}

\end{document}
